\documentclass[10pt,a4paper]{amsart}
\usepackage[margin=19mm]{geometry}
\usepackage{amsmath,amssymb,mathtools}
\usepackage[scr=rsfs,cal=euler]{mathalfa}
\usepackage{xfrac}
\usepackage[unicode,hidelinks,bookmarks=false]{hyperref}
\newcommand{\ie}{i.e.\ }
\newcommand{\cf}{cf.\ }
\newcommand{\resp}{respectively\ }
\newcommand{\Subsection}{\subsection}
\providecommand{\nobreakdash}{\nobreak\hbox{-}}
\setlength{\emergencystretch}{2em}
\multlinegap0pt
\usepackage{amssymb}
\usepackage{tikz}
\usepackage{tcolorbox}
\usepackage[shortlabels]{enumitem}
\newtheorem{theorem}{Theorem}
\title{Sorting from Counterexamples}
\author{Noga Alon, Shay Moran, Shlomo Moran}
\date{Source: arXiv:2608.21579v1 (CC BY 4.0)}
\begin{document}
\maketitle

\noindent\emph{Source and licence.} Sorting from Counterexamples. Version arXiv:2608.21579v1 (CC BY 4.0), distributed by arXiv under the Creative Commons Attribution 4.0 International licence, http://creativecommons.org/licenses/by/4.0/, as recorded in the arXiv licence metadata for this exact version and shown on the abstract page https://arxiv.org/abs/2608.21579v1. Full licence notice: \texttt{licenses/arxiv-2608.21579-CC-BY-4.0.txt}.\par\smallskip

\noindent\textit{Editorial scope.} Counted unit: the article's theorem thm:arbitrary-lower (source0.tex lines 1557--1564). The article's complete original proof of it is delivered with it (source0.tex lines 1566--1754, one \texttt{proof} environment of 189 lines). No mathematics is written, repaired or re-derived: the statement and the proof are the article's own lines, in the article's own notation, and no retained line is substituted or re-flowed.  No further source-local context is needed: every cross-reference of every retained line resolves inside the two delivered spans.  Nothing was omitted from any retained span, so the package carries no omission record.  No retained line contains a citation, so the wrapper carries no bibliography and no bibliography entry.  No retained line contains a figure environment, an image-inclusion command or drawing code, so no image is delivered and the package declares no asset dependency.  Editorial formatting only, in addition: an AMS \texttt{amsart} preamble that restates the article's own declarations and macros used by the retained lines, namely the environments \texttt{theorem} on separate counters, together with the standard packages those lines need.  The retained environments print in this document as Theorem 1 in order of appearance, whereas the article prints the same items with the numbering of its own full document.  The complete native source of the article as distributed in the arXiv e-print for this exact version is bundled unchanged as \texttt{sources/208258/source0.tex}.
\begin{theorem}
\label{thm:arbitrary-lower}
For every $n\geq 3$ and $k\geq 0$, every possibly randomized proper learner
for arbitrary rankings has worst-case expected query complexity
\[
    \Omega(n\log n+nk).
\]
\end{theorem}

\begin{proof}
We prove the two terms separately.

We first establish the $\Omega(n\log n)$ lower bound using only truthful
counterexamples. For simplicity assume that $n$ is a power of two; for
general $n$, we may restrict attention to the largest power of two below
$n$. Consider the comparison tree induced by MergeSort. Every
root-to-leaf path in this tree has length
\[
    D=\Omega(n\log n).
\]
Indeed, at every level of the recursion, merging pairs of equal-size lists
requires altogether $\Omega(n)$ comparisons, and there are
$\Theta(\log n)$ levels.

Choose a random root-to-leaf branch of this tree by choosing each of the two
children with probability $1/2$ at every internal node. The resulting leaf
determines a target ranking $\pi^\star$, which is thus fixed before the
interaction begins.

We describe a truthful oracle for this target. The oracle maintains a
current node on the chosen branch, initially the root. Suppose that the
learner proposes a ranking $\pi$. If $\pi$ disagrees with the outcome of
some comparison on the portion of the branch already traversed, the oracle
may simply return such a comparison as a truthful counterexample.

Otherwise, starting from the current node, the oracle follows the chosen
branch down the comparison tree. As long as $\pi$ agrees with the outcome
of the comparison at the current node, the oracle proceeds to the next
node without returning anything. At the first comparison on which $\pi$
disagrees with the chosen branch, the oracle returns that pair as a
counterexample and makes this node part of the traversed portion of the
branch. If no such comparison is encountered before reaching the leaf,
then $\pi$ agrees with every comparison on the root-to-leaf path and hence
$\pi=\pi^\star$; the oracle then declares the ranking correct.

All counterexamples produced by this strategy are truthful. We now bound
how quickly a learner can move down the random branch. Let $A_t$ denote the
number of previously untraversed edges of the comparison tree that are traversed on
the $t$-th query. If the learner contradicts an earlier comparison then~$A_t=0$. Otherwise, conditional on everything that has happened so far and
on the learner's current query, every new branch choice encountered by the
oracle is still an independent fair choice between the two children. At
each such node, the learner's ranking prescribes one of the two possible
outcomes. Hence it agrees with the random branch with probability $1/2$,
and disagrees with probability $1/2$.

Consequently, $A_t$ is dominated by a geometric random variable with mean
$2$, and therefore
\[
    \mathbb E[A_t\mid\text{history before query }t]\leq 2.
\]
This remains true when the learner is randomized, since we may condition
also on its internal randomness up to this point.

Let $Q$ be the number of queries made before the target is identified. By
the time the interaction terminates, the entire root-to-leaf path has been
traversed, so
\[
    \sum_{t=1}^{Q} A_t\geq D.
\]
Therefore
\begin{align*}
    D
    &\leq
    \mathbb E\left[\sum_{t=1}^{Q}A_t\right] \\
    &=
    \sum_{t\geq1}
    \mathbb E\left[
        \mathbf 1_{\{Q\geq t\}}A_t
    \right] \\
    &\leq
    2\sum_{t\geq1}\Pr(Q\geq t)
    =
    2\mathbb E[Q].
\end{align*}
Thus
\[
    \mathbb E[Q]\geq D/2=\Omega(n\log n),
\]
where the expectation is over both the random branch and the learner's
randomness. Averaging over the choice of the branch, there must therefore
exist a fixed branch, and hence a fixed target ranking $\pi^\star$, for
which the learner requires $\Omega(n\log n)$ queries in expectation.

We next establish the $\Omega(nk)$ term. Consider the directed cycle
\[
    1\to2\to\cdots\to n\to1,
\]
and denote its edges by $e_1,\ldots,e_n$. For each $r\in[n]$, let
$\sigma_r$ be the linear order obtained by deleting $e_r$ from the cycle
and taking the resulting directed path as the order. Thus $\sigma_r$ agrees
with every cycle edge except $e_r$.

Every linear order $\pi$ violates at least one edge of the cycle. Moreover,
if $\pi$ violates $e_r$ and the oracle returns this pair, then the
counterexample is truthful with respect to every $\sigma_s$, $s\neq r$,
and untruthful with respect to $\sigma_r$.

Define a reference interaction that does not depend on the target. Whenever
the learner proposes a ranking $\pi$, choose, according to any fixed rule,
one cycle edge violated by $\pi$ and return it. For example, choose the
violated edge of smallest index. This produces an infinite sequence of
edges
\[
    e_{r_1},e_{r_2},\ldots
\]
that depends on the learner's randomness, but not on the choice of the
target.

For each $r\in[n]$, let $\tau_r$ be the round on which $e_r$ is returned
for the $(k+1)$-st time in this reference interaction, with
$\tau_r=\infty$ if this never occurs. Let
\[
    \tau_{(1)}\leq\tau_{(2)}\leq\cdots\leq\tau_{(n)}
\]
be these values arranged in increasing order. Then
\[
    \tau_{(j)}\geq j(k+1),
    \qquad j=1,\ldots,n.
    \tag{21}\label{eq:cycle-hitting-times}
\]
Indeed, by round $\tau_{(j)}$, at least $j$ distinct cycle edges have each
been returned at least $k+1$ times, so at least $j(k+1)$ counterexamples
must have been returned in total. Consequently,
\[
    \frac1n\sum_{r=1}^n \tau_r
    \geq
    \frac1n\sum_{j=1}^n j(k+1)
    =
    \frac{(n+1)(k+1)}2.
    \tag{22}\label{eq:average-cycle-hitting-time}
\]
If some $\tau_r=\infty$, this inequality is immediate.

Now choose $R$ uniformly from $[n]$, independently of the learner, and set
\[
    \pi^\star=\sigma_R.
\]
Thus the target is fixed before the interaction begins.

We define an oracle for this target by following the reference interaction
until round $\tau_R$. Before this round, the edge $e_R$ has been returned
at most $k$ times. Each occurrence of $e_R$ is an untruthful counterexample
with respect to $\sigma_R$, while every other returned cycle edge is
truthful. Hence all these responses are valid and use at most $k$
untruthful counterexamples.

On round $\tau_R$, the reference rule is about to return $e_R$ for the
$(k+1)$-st time. If the learner's query is $\sigma_R$, the oracle declares
it correct. Otherwise, the query violates some cycle edge other than
$e_R$, and the oracle returns such an edge instead; this counterexample is
truthful. Indeed, a ranking whose only violated cycle edge is $e_R$ must be
exactly $\sigma_R$. From this point on, whenever the learner proposes a
ranking different from $\sigma_R$, the oracle may return any truthful
counterexample, and it declares the ranking correct when $\sigma_R$ is
proposed. Thus this defines a valid oracle using at most $k$ untruthful
counterexamples.

Let $Q$ denote the number of queries before the interaction terminates.
Under the coupling above, the learner cannot terminate before round
$\tau_R$, and therefore
\[
    Q\geq\tau_R.
\]
Taking expectation over both the random choice of $R$ and the learner's
internal randomness, and using \eqref{eq:average-cycle-hitting-time}, gives
\[
\begin{aligned}
    \mathbb E[Q]
    &\geq \mathbb E[\tau_R] \\
    &= \mathbb E\left[\frac1n\sum_{r=1}^n\tau_r\right] \\
    &\geq \frac{(n+1)(k+1)}2
    = \Omega(nk).
\end{aligned}
\]
Averaging over $R$, there is therefore some fixed target $\sigma_R$ for
which the expected number of queries, over the learner's randomness, is
$\Omega(nk)$.

We have proved that the worst-case expected query complexity of every
randomized proper learner is both $\Omega(n\log n)$ and $\Omega(nk)$.
Consequently,
\[
    \max\{\Omega(n\log n),\Omega(nk)\}
    =
    \Omega(n\log n+nk),
\]
up to a universal constant.
\end{proof}

\end{document}
